\section*{Chapter \ref{ch:b2:boltzmann-factor} --- The Boltzmann Factor}

\begin{solution}{exo:b2:boltzmann-factor:1}
$H = k_BT/mg = \qty{8.4}{km}$; \qty{0.70}{}, \qty{0.35}{}, \qty{0.30}{bar} (the real
\qty{0.26}{}: the air cools with height); helium \qty{61}{km}; Mars
\qty{10.7}{km}.
\end{solution}

\begin{solution}{exo:b2:boltzmann-factor:2}
Hydrogen: $4\eu^{-395} \approx 10^{-171}$; $4\eu^{-19.7} = 10^{-8}$; $4\eu^{-11.8} = 3 \times
10^{-5}$. Sodium: $\eu^{-9.8} = 6 \times 10^{-5}$. Rotation: $\eu^{-0.04} = 0.96$, almost
equal. Nuclear spin: $1 - 4 \times 10^{-6}$ --- the tiny polarisation that
magnetic resonance works with.
\end{solution}

\begin{solution}{exo:b2:boltzmann-factor:3}
N$_2$ \qty{300}{K}: \qty{422}{}, \qty{476}{}, \qty{517}{m/s}; \qty{77}{K}: \qty{214}{}, \qty{241}{},
\qty{262}{m/s}; H$_2$: \qty{1580}{}, \qty{1780}{}, \qty{1930}{m/s}; smoke particle:
\qty{9}{cm/s} (its Brownian shiver). Above $2v_{\text{p}}$ one molecule in
twenty; above $3v_{\text{p}}$ one in $2500$ --- still $10^{22}$ per cubic metre.
\end{solution}

\begin{solution}{exo:b2:boltzmann-factor:4}
$x = 2.2 \times 10^{-3}$ and $0.67$; $M/M_{\text{sat}} = 2.2 \times 10^{-3}$ and $0.59$;
$M_{\text{sat}} = n\mu = \qty{9.3e4}{A/m}$; $\chi = \mu_0n\mu^2/k_BT = 2.6 \times 10^{-4}$.
Iron's \qty{1.7e6}{A/m} is cooperative, not a Boltzmann effect.
\end{solution}

\begin{solution}{exo:b2:boltzmann-factor:5}
(a) \cref{prop:b2:boltzmann-factor:twolevel}. (b) $\eu^{-x} < 1$ for $T > 0$;
$N_2 > N_1$ would need $T < 0$ --- a \omterm{prop:b2:laser:gain}{population inversion}. (c) $0$ and
$Nk_B\ln 2$. (d) $x = 1.93$: $0.13$; at \qty{77}{K}, $x = 7.5$: $5 \times 10^{-4}$.
\end{solution}

\begin{solution}{exo:b2:boltzmann-factor:6}
(a) Derivative of $N\varepsilon/(\eu^x + 1)$. (b) $x = 2.40$, $C = 0.44\,Nk_B$.
(c) $\propto x^2\eu^{-x} \to 0$: frozen; $\propto 1/T^2 \to 0$: both levels full,
nothing left to absorb. (d) $k_BT = 0.42\varepsilon$: \qty{0.5}{mK}; below a
millikelvin the nuclear spins hold most of the heat capacity and
slow every cooling.
\end{solution}

\begin{solution}{exo:b2:boltzmann-factor:7}
(a) \cref{thm:b2:boltzmann-factor:equipartition}. (b) $\tfrac32R$,
$\tfrac52R$. (c) $C_{\text{vib}}/R = x^2\eu^x/(\eu^x - 1)^2$: $0.002$, $0.41$, $0.90$;
$C_V = 2.50R$, $2.91R$, $3.40R$. (d) $0.92$, $0.30$, $0.995$; diamond $x = 4.3$:
$0.26 \times 3R = \qty{6.5}{J/mol/K}$.
\end{solution}

\begin{solution}{exo:b2:boltzmann-factor:8}
(a), (b) \cref{thm:b2:boltzmann-factor:maxwell}. (c) $\tfrac32k_BT$. (d)
$\tfrac14n\langle v\rangle = \qty{2.9e27}{m^{-2}.s^{-1}}$: $3 \times 10^{23}$ per square
centimetre per second; through \qty{7.9e-13}{m^2}: $2.3 \times 10^{15}$
molecules per second, a leak of $\qty{1e-5}{Pa.m^3/s}$.
\end{solution}

\begin{solution}{exo:b2:boltzmann-factor:9}
(a) Integrate $f$ by parts; the leading term is $(2/\sqrt\pi)x\eu^{-x^2}$. (b)
Helium $x = 5.5$: $5 \times 10^{-13}$; nitrogen $x = 14.5$: $\eu^{-211}$. (c) $x =
4.9$: $10^{-10}$ per collision time --- gone; $10^{-17}$ needs $x \approx 6.6$,
$T \lesssim \qty{230}{K}$. (d) $x = 11$: $\eu^{-121}$; cold enough.
\end{solution}

\begin{solution}{exo:b2:boltzmann-factor:10}
(a) $n = n_0\eu^{m\omega^2r^2/2k_BT}$. (b) $\eu^{\Delta m\omega^2r^2/2k_BT} = \eu^{0.24} =
1.27$. (c) $H = k_BT/m'g = \qty{16}{\micro m}$; counting spheres at
several heights gives $k_B$, hence $N_A = R/k_B$. (d) Ions follow
$\eu^{\mp e\varphi/k_BT}$ around a charge; linearised in Poisson's equation
this gives $\varphi'' = \varphi/\lambda_{\text{D}}^2$: the potential is screened
beyond $\lambda_{\text{D}}$.
\end{solution}

\begin{solution}{exo:b2:boltzmann-factor:11}
(a) $Z = 1/(1 - \eu^{-x})$, $\langle E\rangle = h\nu/(\eu^x - 1)$. (b) $k_BT$:
equipartition; $h\nu\eu^{-h\nu/k_BT}$: frozen. (c) $u_\nu = (8\pi\nu^2/c^3)\langle E
\rangle$; Rayleigh--Jeans gives each mode $k_BT$. (d) A mode cannot be
excited by less than one quantum.
\end{solution}

\begin{solution}{exo:b2:boltzmann-factor:12}
(a) $x = \mu B/k_BT$ only; $\ln 2$ and $0$. (b) $x = 0.67$: $S/Nk_B = 0.51$, so
$0.18\,k_B$ per spin removed. (c) $B/T$ constant: \qty{10}{mK}. (d) The spins'
own field (millitesla) replaces $B$ at the end and caps the cooling;
measuring $M = n\mu\tanh(\mu B/k_BT)$ gives $T$.
\end{solution}

\begin{solution}{pb:b2:boltzmann-factor:1}
\textbf{1.} $\dd P/\dd z = -nmg = -(mg/k_BT)P$; $H = \qty{8.4}{km}$.

\textbf{2.} $\eu^{-mgz/k_BT}$: the potential energy of one molecule, the
air's temperature.

\textbf{3.} \qty{0.70}{}, \qty{0.52}{}, \qty{0.35}{bar}; $H\ln 2 = \qty{5.8}{km}$.

\textbf{4.} $P_0/g = \qty{1.03e4}{kg/m^2}$; $\times 4\pi R^2$: \qty{5.3e18}{kg}; $\rho_0H =
1.22 \times 8400 = \qty{1.03e4}{kg/m^2}$.

\textbf{5.} $5.3 \times 10^{18}/(29 \times 10^{-3}/N_A) = 1.1 \times 10^{44}$.

\textbf{6.} $H_{\text{O}_2} = 7.6$, $H_{\text{N}_2} = \qty{8.7}{km}$: the ratio at \qty{50}{km}
would be $\eu^{-50/7.6 + 50/8.7} = 0.44$ of its ground value; turbulence mixes
faster than diffusion separates.

\textbf{7.} Lower: colder air is denser and the pressure falls faster.

\textbf{8.} \qty{2.5e25}{m^{-3}}; $\eu^{-11.9}$: \qty{1.7e20}{m^{-3}} --- no edge, only a
convention.

\textbf{9.} $11.2\sqrt{6370/6870} = \qty{10.8}{km/s}$.

\textbf{10.} Helium \qty{2040}{}, \qty{2300}{m/s}; nitrogen \qty{770}{}, \qty{870}{m/s}.

\textbf{11.} $x = 5.3$: $4 \times 10^{-12}$; nitrogen $\eu^{-196}$: none.

\textbf{12.} $\tfrac14 \times 10^{12} \times 2300 \times 4 \times 10^{-12} \times \tfrac12 \approx
10^3$ atoms per square metre per second; over the Earth and a year,
$2 \times 10^{25}$ atoms --- \qty{0.1}{kg}.

\textbf{13.} $5 \times 10^{-6} \times 1.1 \times 10^{44} = 5.5 \times 10^{38}$ atoms, \qty{3.7e12}{kg};
against \qty{0.1}{kg/yr} the residence time would be absurd.

\textbf{14.} Steady state requires a loss of \qty{3e6}{kg/yr}: a residence
time of about a million years, which is the accepted value --- so the
thermal estimate at \qty{1000}{K} is far too small: the escape is
dominated by hotter episodes and by non-thermal (ionic) processes.

\textbf{15.} $x = 3.75$: $\eu^{-14}$ instead of $\eu^{-28}$ --- a factor $10^6$:
the loss is extraordinarily sensitive to the exospheric temperature.

\textbf{16.} The Moon is hot and light: everything escapes in
geological time; Titan is cold, $x = 11$ for nitrogen.

\textbf{17.} $p_\pm = \eu^{\pm x}/2\cosh x$; $M = n\mu\tanh x$.

\textbf{18.} $x = \mu B/k_BT$: \qty{4.2}{A/m}, \qty{310}{A/m}, \qty{1.6e4}{A/m}
(saturation \qty{1.9e4}{A/m}); Curie for $x \ll 1$, i.e. $T \gg \qty{70}{mK}$.

\textbf{19.} $M$ depends on $T$ alone once $B$ is known; steepest near
$x \approx 1$, $T \approx \mu B/k_B = \qty{70}{mK}$; one measures the
susceptibility with a coil.

\textbf{20.} $\langle E\rangle = -N\mu B\tanh x$; $C = Nk_Bx^2/\cosh^2x$, peaking at
$x = 1.2$: \qty{56}{mK}.

\textbf{21.} $Nk_B\ln 2$ and $0$: all the spins align with the field.

\textbf{22.} \qty{10}{mK}; the spins' own field of a millitesla replaces $B$
and fixes the floor.

\textbf{23.} $x = 1$ at $\mu_{\text{N}}B/k_B = \qty{0.37}{mK}$: Curie holds down to a
millikelvin, and the thermometer is useful to microkelvins.

\textbf{24.} $\Delta\nu/\nu = \sqrt{8k_BT\ln2/mc^2} \propto \sqrt T$: the width
of a line measures the gas temperature.

\textbf{25.} Atmosphere, escaping gas, spins: $\eu^{-E/k_BT}$ with $E = mgz$,
$\tfrac12mv^2$, $\mp\mu B$.
\end{solution}
